\section*{第 \ref{ch:g10:measuring-the-earth} 章——测量地球}

\begin{solution}{exo:g10:measuring-the-earth:1}
船只消失时先没入船身：若海面是平的，
船只只会越来越小，全部船身仍可见。月食时地球在月球上的影子总是圆弧，
若地球是平圆盘，侧看时影子会拉长成细线。向北行进时，北极星会升至地平线以上，
若地球是平的，北极星的高度
应几乎不变。
\end{solution}

\begin{solution}{exo:g10:measuring-the-earth:2}
\omterm{def:g10:measuring-the-earth:geodesy}{大地测量学}测定地球的形状和大小，以及地面各点的位置，
并研究它们随时间的变化。\omterm{def:g10:measuring-the-earth:triangulation}{三角测量}通过一串三角形定位各点，
其中所有角都要测量；地面上只需测量一条长度，
即基线。
\end{solution}

\begin{solution}{exo:g10:measuring-the-earth:3}
$C = 50 \times 5000 = 250\,000$ 斯塔德；$250\,000 \times \qty{157.2}{m} =
\qty{39300}{km}$。
\end{solution}

\begin{solution}{exo:g10:measuring-the-earth:4}
$111/60 \approx \qty{1.85}{km}$。四分之一经线：$90 \times 111 \approx
\qty{10000}{km}$。
\end{solution}

\begin{solution}{exo:g10:measuring-the-earth:5}
$a - b = 6378.137 - 6356.752 = \qty{21.385}{km}$；
$f = 21.385/6378.137 \approx 0.003353 \approx 1/298$。
\end{solution}

\begin{solution}{exo:g10:measuring-the-earth:6}
南方：$\tan\theta_1 = 0.364$，$\theta_1 \approx \qty{20.0}{\degree}$。北方：
$\tan\theta_2 = 0.456$，$\theta_2 \approx \qty{24.5}{\degree}$。
$\alpha = \qty{4.5}{\degree}$，因此 $C = (360/4.5) \times 500 = 80 \times 500 =
\qty{40000}{km}$。
\end{solution}

\begin{solution}{exo:g10:measuring-the-earth:7}
$C = 180 - 62 - 71 = \qty{47}{\degree}$。
$AC = 10.00 \sin 71^\circ/\sin 47^\circ = 10.00 \times 0.9455/0.7314 \approx
\qty{12.93}{km}$；$BC = 10.00 \sin 62^\circ/\sin 47^\circ \approx
\qty{12.07}{km}$。
\end{solution}

\begin{solution}{exo:g10:measuring-the-earth:8}
经度的变化为 $10 + 20 = \qty{30}{\degree}$。在
北纬 \qty{60}{\degree}，一度的长度为 $111.3 \times 0.5 = \qty{55.65}{km}$，因此
距离为 $30 \times 55.65 \approx \qty{1670}{km}$。在赤道上：
$30 \times 111.3 \approx \qty{3340}{km}$，是前者的两倍。
\end{solution}

\begin{solution}{exo:g10:measuring-the-earth:9}
视线与海面相切于 $T$ 点：圆的切线
垂直于过切点的半径，因此 $T$ 处的角为
直角。由毕达哥拉斯定理，$(R+h)^2 = R^2 + d^2$，故 $d^2 = 2Rh + h^2$；$h$ 与
$R$ 相比很小，所以 $d \approx \sqrt{2Rh} = \sqrt{2 \times
6\,371\,000 \times 2} \approx \qty{5050}{m}$，约为 \qty{5}{km}。
\end{solution}

\begin{solution}{exo:g10:measuring-the-earth:10}
曲线在中纬度地区最陡，约在 \qty{30}{\degree} 至
\qty{60}{\degree} 之间，在赤道和极点附近则较平缓。在扁球形的地球上，
极地附近的表面弯曲程度较小；表面越平缓，
竖直方向（从而纬度）变化一度，
所需走过的距离就越远。
\end{solution}

\begin{solution}{exo:g10:measuring-the-earth:11}
$1072.2/9.646 \approx \qty{111.15}{km}$ 每度；$90 \times 111.15
\approx \qty{10004}{km}$，比真实的四分之一经线多 \qty{2}{km}。
弧段中心位于 \qty{46}{\degree} 附近，那里一度的长度略大于
从赤道到极点的平均一度
（$10\,002.0/90 \approx \qty{111.13}{km}$）。
\end{solution}

\begin{solution}{exo:g10:measuring-the-earth:12}
取 \qty{7.1}{\degree}：$C = (360/7.1) \times 5000 \approx 253\,500$ 斯塔德
而非 $250\,000$：约大了 1.4\,\%（同样的距离对应较小的角，
意味着地球更大）。沿尼罗河由步行者量出的距离，
远不如测得的角度精确。
\end{solution}

\begin{solution}{exo:g10:measuring-the-earth:13}
在 $ABC$ 中：$C = \qty{50}{\degree}$，$AC = 6.000 \sin 75^\circ/\sin 50^\circ
\approx \qty{7.566}{km}$。在 $ACD$ 中：$D$ 处的角为 $180 - 50 - 68 =
\qty{62}{\degree}$，$AD = 7.566 \sin 68^\circ/\sin 62^\circ \approx
\qty{7.945}{km}$。
\end{solution}

\begin{solution}{exo:g10:measuring-the-earth:14}
$\cos\varphi = 1/2$ 得 $\varphi = \qty{60}{\degree}$；$\cos\varphi = 0.1$
得 $\varphi \approx \qty{84.3}{\degree}$。纬度 $\varphi$ 处的经度一度，
其长度等于赤道上一度的长度乘以 $\cos\varphi$：
该比值只取决于角度，与半径无关。
\end{solution}

\begin{solution}{exo:g10:measuring-the-earth:15}
在赤道：$111.195 - 110.574 = \qty{0.621}{km} = \qty{621}{m}$；在
极点：$111.694 - 111.195 = \qty{0.499}{km} = \qty{499}{m}$。每度误差达数百米，
逐度累加后远大于一米，
由纬度和经度换算位置时，必须采用
椭球体。
\end{solution}

\begin{solution}{pb:g10:measuring-the-earth:1}
\textbf{1.} $31.20 - 24.09 = \qty{7.11}{\degree}$。

\textbf{2.} $7.11/360 \approx 0.0198 = 1/50.6$，非常接近
埃拉托色尼的 $1/50$。

\textbf{3.} $C = 788 \times 360/7.11 \approx \qty{39900}{km}$。

\textbf{4.} 两纬线之间的子午线弧：正午角之差
即纬度之差，它只表示两地的南北间距，
两地距离中的东西向部分
不起作用。

\textbf{5.} 地方正午时太阳位于子午面内，两地皆然；同一天它在恒星间的位置相同，
因此两个角度的差异
仅源于地球的曲率。

\textbf{6.} $180 - 58 - 66 = \qty{56}{\degree}$。

\textbf{7.} $AC = 12.000 \sin 66^\circ/\sin 56^\circ = 12.000 \times
0.9135/0.8290 \approx \qty{13.223}{km}$。

\textbf{8.} $BC = 12.000 \sin 58^\circ/\sin 56^\circ \approx \qty{12.275}{km}$。

\textbf{9.} 点 $D$ 处的角为 $180 - 47 - 75 = \qty{58}{\degree}$；
$AD = 13.223 \sin 75^\circ/\sin 58^\circ \approx \qty{15.061}{km}$。

\textbf{10.} $15.061 \cos 8^\circ \approx \qty{14.914}{km}$。

\textbf{11.} $14.914/111.13 \approx \qty{0.1342}{\degree}$，即
$0.1342 \times 60 \approx 8.05$ 角分。

\textbf{12.} $111.51 - 110.57 = \qty{0.94}{km}$：\qty{940}{m}。

\textbf{13.} 靠近极点时，竖直方向随距离转动得
更慢：那里的地表弯曲较小。地球在两极处扁平，
在赤道处隆起。

\textbf{14.} $a - b = \qty{21.385}{km}$；$f = 21.385/6378.137 \approx 1/298$。

\textbf{15.} 铅垂线垂直于椭球体的表面。在椭圆上，
曲线的法线仅在赤道和两极处经过中心；
在 \qty{45}{\degree} 处，它与轴的交点
位于中心的下方。

\textbf{16.} $10\,000\,000$ 米，即 \qty{10000}{km}。

\textbf{17.} $10\,001\,966/10\,000\,000 \approx \qty{1.0001966}{m}$。

\textbf{18.} 偏短：真实的千万分之一是 1799 米的
$1.0001966$ 倍，因此 1799 米约短了其长度的 $1.97 \times 10^{-4}$，
即 0.02\,\%。

\textbf{19.} 偏短。该标尺取自 1790 年代计算的
四分之一子午线的千万分之一；标尺既短于
真实四分之一子午线的千万分之一，
故计算所得的四分之一子午线比真值短。

\textbf{20.} $1.0001966 - 1 = 0.0001966$ 米：第一米约
\textbf{偏短 \qty{0.2}{mm}}——相当于一米内约两张纸的厚度，
这一精度对于借助望远镜、跨越整个国家的角度测量来说，
实在令人惊叹。
\end{solution}
